Emmy was nog op pad.
Al dolend door de stad
en met teelsap beklad,

werd zij en de jachtstoet
door een priester begroet
bij een statig landgoed.

Een prestigieuze kerk
op deze vrome perk
was het geprezen werk

van grootste devotie.
Hier strandde voor Emmy
de vuige processie.

Ondanks de bedreiging
was de vloerverwarming
een fijne ervaring.

De jager boog aldaar
voor een antiek altaar
en benagelde haar

in een grote kerkzaal
aan een eiken schandpaal
en nam haar andermaal.

Haar handen waren vrij
als hulp bij het gevrij
en maakten de weg vrij.

Hij greep haar in haar zij
en trok zo elke dij
voldoende dichterbij.

Dit keer was er publiek.
Die voorzag zijn techniek
van bruikbare kritiek

en opjuttend gejuich.
Het schijnheilige tuig
maakte hem extra ruig

in het brandje blussen
binnen haar stootkussen.
Zij kon ondertussen

bij zijn amuseren
haar brein penetreren
om te reflecteren;

“Is dit mijn faire straf?
Omdat ik er zo laf
ineens de brui aan gaf.

Als zonden verdorden
als er genoeg horden
koen genomen worden,

hoeveel moet ik dan nog?
Als ik mijn plagen log,
dan zijn er genoeg toch?

Ik ben immers ontmaagd.
Heb geen minuut geklaagd.
Ik ben daarop verlaagd.

Niet slechts in het geniep!
Terwijl ik naakt rondliep
en men mij grof toeriep,

hebben ze valselijk
en nogal schadelijk
mijn stand publiekelijk

heel diepgaand vernederd.
Mijn aanzien is sedert
nimmer meer gevederd.”

De kerk leek omgebouwd
alsof men welbeschouwd
het oude had mistrouwd.

De nieuwe kerk was groot.
Ramen met glas in lood.
Niet over Jezus’ dood

verhaalden de prenten,
maar van incidenten,
waarin delinquenten

diep waren gevallen.
Vrouwen die zich stallen
aan het lid van allen.

De Samaritaanse
die al met vijf sjanste.
Met een man herkanste

die niet haar gemaal was.
Die lompe liefdespas was
voor de Messias

geen veroordeling waard.
Hij gaf de vrouw bedaard
water dat de ziel spaart.

Of de trouwgewonde,
die men wierp ten gronde.
“Wie vrij is van zonde

werpe de eerste steen.”
was wat door het raam scheen.
Direct daarnaast verscheen

een hoer die moest boeten.
Zalvend Jezus’ voeten
kon de heer haar groeten

als lid van zijn stamboom.
Tamar was ook niet vroom.
Rachad had ook geen schroom.

Naast deze panelen
werden taferelen
mede met penselen

of met beitels verbeeld.
En het Mariabeeld
werd alhier uitgebeeld

via Magdalena
als ware Maria
en niet Jezus’ mama.

De klokken begonnen.
De kerkse pionnen
werden zo gewonnen

tot Godshuis was gevuld.
De priester was gehuld
in engelengeduld;

de annunciatie
in een illustratie
op zijn grijze kloffie.

Daar werd de mis een daad
van het schijnheilig kwaad
Hij pleegde een misdaad,

maar deed het geduldig.
Hij groette veelvuldig
tot de laatste schuldig

zijn excuses aanbood
voor zijn hoge plasnood.
Met deze dorpsgenoot

was de sekte compleet.
De priester was gereed
toen men de deur dichtdeed.

“Sodom en Gomorra,
in schrift op een nota,
is het anathema

dat we van de buren
in moeilijke uren
vaak moeten verduren.

En dat geldt ook nu weer.
Abraham vroeg de heer
om gratie keer op keer,

mits er enkelen zijn
die men kan zien als rein
volgens God zijn richtlijn.

God zond twee engelen
die moesten hengelen,
door aan te zwengelen,

met wellust als vislijn,
of er ook burgers zijn,
al was het aantal klein,

waardig om te sparen.
Het tuig wilde paren
en kennis vergaren

als ongepaste lust.
Deze hete onrust
kreeg Lot zelfs niet geblust,

toen hij uit hoge nood
zijn twee dochters aanbood
in hun maagdelijk bloot.

Men weigerde dat bod.
Daarom ontsnapte Lot
en vernietigde God

de goddeloze stad.
Waaruit we leren dat
het weigeren van wat

aan ons wordt gegeven,
met ons eigen leven
ooit wordt afgeschreven.

De jager biedt Emmy
in ruil voor clementie.
Ik geef zijn vrouw gratie

en neem deze ruilvrouw,
als teken van zijn trouw;
als teken van berouw.

Onderdruk uw aandrift,
trek terug uw klaagschrift
en zie dit als een gift.

Weer schonk onze broeder
een kersverse loeder.
God is de behoeder.”

Emmy kreeg deze preek,
maar amper mee en leek
voornamelijk van streek

door de vuistgebaren.
Hoon had ze ervaren.
Fruit zat in haar haren.

Haar lijf was afgekneld.
Werd ze ook gegeseld
en aan de kaak gesteld

na het publiek minnen?
Emmy schoot te binnen
de paap die zijn zinnen

over plaatsen begon
die je als prille non,
maar beter mijden kon.

Daar kende men schoften,
bizarre bruiloften,
huwelijksgeloften

die hij verafschuwde.
Dat de vader spuwde
over de gehuwde,

haar borsten en haar hoofd,
omdat men daar geloofd
dat God haar dan beloofd

dat ze niet wordt versteend.
Ook werd er daar gemeend
dat God zijn steun verleent

aan het verse bruidspaar
als hun kleding en haar
met rottend etenswaar,

zoals bloem of een ei,
maar tevens allerlei
ander smerig geklei,

door de clan werd versierd.
Was het ongemanierd
dat men het trouwen viert

door een heidens gebruik?
Emmy haar onderbuik
voelde als een melkkruik.

Ze overwoog toen dat
in deze vreemde stad
men dit ritueel had

voor elke trouwerij.
Deze bruiloftspartij
hoorde er gewoon bij

als lokale gekte.
Haar man die haar dekte
was lid van een sekte

voor wie dit gewoon was.
Emmy vermoedend las
de jager zijn grimas

en aldus voorspelde,
ondanks haar ellende,
het stilaan bekende.

Terwijl hij in zijn gast
zijn witte vruchten plast,
zat haar nek stevig vast.

Ze keek zo bereden
beduusd naar beneden
waar gekleurde treden

leiden naar het altaar.
De jager greep haar haar
en kwam daarna klaar.

De kerkvloer van tegels
beschilderd met regels
bedacht door de vlegels,

die haar vernederden
en haar herinnerden
aan vele honderden

verwante voorgangers,
werd door de kerkgangers
gebruikt om handlangers

voor de kerk te winnen
middels krasse zinnen
over het beminnen.

Zo stond er geschreven;
“Preutsheid wordt verdreven.”
en “Zaad is het streven.”

Of  “De vrouw dient de man”,
“Seks is ons pakkie-an”,
“Een vrouw is in de ban

nopens de communie
behalve als hostie”,
“Alles is religie

wat gebeurt in de kerk”,
“De priester doet God’s werk,
zijn celibaat is sterk”,

“Elke biecht is veilig”,
“Elk offer is heilig”
en “Een gift is geilig”.

Al was de jager klaar,
dat gold nog niet voor haar,
geboeid aan het altaar.

De zaadlozende mens
ondervond zijn sekswens
immers bij haar intens.

Hij trok zich dus terug
en stak vervolgens vlug
tussen armen en rug

in lijn met het schandblok
zijn houten wandelstok.
Ze was nu echt de bok.
